% !TEX TS-program = pdflatex
% !TeX encoding = UTF-8
% !TeX spellcheck = en_GB

\documentclass[aspectratio=43]{beamer}
% use this instead for 16:9 aspect ratio:
%\documentclass[aspectratio=169]{beamer}
% supported acpect ratios  1610  169 149 54 43 (default) 32
%

\usepackage[english]{babel} 
\usepackage[utf8]{inputenc}
\usepackage[T1]{fontenc}
\usepackage{algorithm,algorithmic}
\usepackage{hyperref}
\usepackage{multimedia}
\usetheme{ETHbeamer}
\useinnertheme{rectangles}

\colorlet{ETHcolor1}{ETHa}
\colorlet{ETHcolor2}{ETHa}

\author{Beat Hubmann}

\title{Seminar in Physics for CSE\\
    Ardizzone et al.:\\	   
	Analyzing Inverse Problems with Invertible Neural Networks}

\date{December 12, 2019}

% uncomment if you do not want to use a department logo
\deplogofalse

% to get section outlines
\AtBeginSection[]
{
  \begin{frame}
    \frametitle{Outline}
    \tableofcontents[currentsection,hideothersubsections]
  \end{frame}
}

\AtBeginSubsection[]
{
  \begin{frame}
    \frametitle{Outline}
    \tableofcontents[currentsection,currentsubsection]
  \end{frame}
}


% ===============
\begin{document}

\titleframe

% \begin{frame}
% 	\frametitle{Contents}
% 	\tableofcontents
% \end{frame}



\section{Motivation and Selected Theory}

\begin{frame}
    
	\frametitle{What Is an Inverse Problem?}
	\begin{center}	
		\begin{tikzpicture}
			\node<1>[anchor=south west, inner sep=0] (image) at (0,0) {\includegraphics[height=5cm]{INN-forward-process.png}};
			\node<2->[anchor=south west, inner sep=0] (image) at (0,0) {\includegraphics[height=5cm]{INN-inverse-process-direct.png}};
			\node<3> [align=center, red, font={\Huge}] at (image.center) {ambiguous};
			% \node<3->[anchor=south west, inner sep=0] (image) at (0,0) {\includegraphics[height=5cm]{inverse_problems2.png}};
			\node<4> [align=center, red, font={\Huge}] at (image.center) {wanted: $p(x|y)$};
		\end{tikzpicture}\cite{Ardizzone2_2019}
	\end{center}
\end{frame}


% \begin{frame}
% 	\frametitle{Then Why Not Just Train a NN for the Inverse?}
% 	\begin{columns}[T]
% 		\column{5cm}
% 			\includegraphics[width=5cm]{INN-inverse-process-direct.png}

% 		\column{5cm}
% 			direct supervised learning of $y \rightarrow x$ problematic:\\ \alert{ambiguous}

% 	\end{columns}
% \end{frame}


\begin{frame}
	\frametitle{Ambiguous Inverse Problem Go-To Nr. 1: MCMC}
	\begin{columns}[T]
		\column{5cm}
			\begin{center}
				\includegraphics{mcmc.jpg}
			\end{center}
		\column{5cm}
			\begin{itemize}
				\item<1->e.g. Metropolis-Hastings algorithm
				\item<1->\alert{expensive}
			\end{itemize}
	\end{columns}
\end{frame}


\begin{frame}
	\frametitle{Inverse Problem Go-To Nr. 2:\\Approximate Bayesian Computation}
	\begin{center}	
		\begin{tikzpicture}
			\node<1>[align=center, black, font={\Huge}] at (image.center) {$\pi(\theta|y) \propto p(y|\theta)\pi(\theta)$};
			\node<2->[anchor=south west, inner sep=0] (image) at (0,0) {\includegraphics[height=6cm]{ABC_parameter_est.png}\cite{Sunnaker2013}};
			\node<3->[align=center, red, font={\Huge}] at (image.center) {expensive};
		\end{tikzpicture}
	\end{center}
\end{frame}

\begin{frame}
    \begin{center}
        \begin{tikzpicture}
            \node<1>[anchor=south west, inner sep=0] (image) at (0,0) {\includegraphics[height=9cm]{ABC.png}\cite{Sunnaker2013}};
            \node<2>[anchor=south west, inner sep=0] (image) at (0,0) {\includegraphics[width=11cm]{ABC1.png}\cite{Sunnaker2013}};
            \node<3>[anchor=south west, inner sep=0] (image) at (0,0) {\includegraphics[width=11cm]{ABC2.png}\cite{Sunnaker2013}};
            \node<4>[anchor=south west, inner sep=0] (image) at (0,0) {\includegraphics[width=11cm]{ABC3.png}\cite{Sunnaker2013}};
		\end{tikzpicture}
    \end{center}
\end{frame}








\begin{frame}
	\frametitle{Neural Network-based Approaches (1)}
	\begin{itemize}
		\item predict fitting parameter of a distribution: \alert{restrictive}
		\begin{center}
			\includegraphics[height=6cm]{predict_fit.png}\cite{Nix1994}
		\end{center}
	\end{itemize}
	
\end{frame}

\begin{frame}
	\frametitle{Neural Network-based Approaches (2)}
	\begin{itemize}
		\item use variational network weights: \alert{still restrictive}
		\begin{center}
			\includegraphics[width=10cm]{uncertain_weights.png}\cite{Blundell2015}
		\end{center}
	\end{itemize}
\end{frame}

\begin{frame}
	\frametitle{Neural Network-based Approaches (3)}
	\begin{itemize}
		\item use conditional Generative Adversarial Networks: \alert{yes, but ...}
		\begin{center}
			\includegraphics[height=6cm]{conditional_GAN.png}\cite{Mirza2014}
		\end{center}
	\end{itemize}
\end{frame}


\begin{frame}
    \frametitle{What is the difference between standard NN and invertible NN?}
    \begin{center}
        \includegraphics[width=11cm]{NNvsINN.png}\cite{Ardizzone1_2019}
    \end{center}
\end{frame}
        








\begin{frame}
	\frametitle{Creating a Bijective Mapping: $x \leftrightarrow [y, z]$}
	\begin{center}
		\begin{tikzpicture}
			\node[anchor=south west, inner sep=0] (image) at (0,0) {\includegraphics[width=10cm]{INN-forward-process-with-z.png}};
		\end{tikzpicture}\cite{Ardizzone2_2019}
	\end{center}
\end{frame}


% \begin{frame}
% 	\frametitle{Invertible Neural Networks to the Rescue!}
% 	\begin{center}
% 		\begin{tikzpicture}
% 			\node[anchor=south west, inner sep=0] (image) at (0,0) {\includegraphics[width=10cm]{INN-training-scheme.png}};
% 		\end{tikzpicture}
% 	\end{center}
% \end{frame}





\begin{frame}
	\frametitle{Resolving the Ambiguity}
	\begin{center}	
		\begin{tikzpicture}
			\node<1-> [anchor=south west, inner sep=0] (image) at (0,0) {\includegraphics[width=10cm]{INN-scheme.png}\cite{Ardizzone2_2019}};
			\draw<1-> [red, thick] (9,0.4) rectangle (10,2.6);
			\node<1> [align=center, white, font={\Large}] at (5,-0.75) {$p(x|y)$ replaced by deterministic function:\\ $x=g(y,z;\theta);\quad z\sim \mathcal{N}(0,1)$};
			\node<1-2> [align=center, white, font={\Large}] at (5,-2) {$\hat{p}(x=g(y,z;\theta)|y) = p_z {\lvert{\det{\frac{\partial g(y,z;\theta)}{\partial [y,z]}}}\rvert}^{-1}$}; 
			\node<2-> [align=center, black, font={\Large}] at (5,-0.75) {$p(x|y)$ replaced by deterministic function:\\ $x=g(y,z;\theta);\quad z\sim \mathcal{N}(0,1)$};
			\node<3-> [align=center, black, font={\Large}] at (5,-2) {$\hat{p}(x|y)=q(x=g(y,z;\theta)|y) = p_z {\lvert{\det{\frac{\partial g(y,z;\theta)}{\partial [y,z]}}}\rvert}^{-1}$}; 
		\end{tikzpicture}
	\end{center}
\end{frame}








\begin{frame}
	\frametitle{Invertible Neural Networks: Core Idea}
	% \begin{center}
			\begin{description}
			\item<1->$s_1, s_2, t_1, t_2: \mathbb{R} \rightarrow \mathbb{R} $
			\item<1->$u_1, u_2 \in \mathbb{R}$
			\item<2->$v_1 = u_1 \cdot \exp(s_2(u_2)) + t_2(u_2)$
			\item<2->$v_2 = u_2 \cdot \exp(s_1(v_1)) + t_1(v_1)$
			\item<3->$u_2 = \frac{v_2 - t_1(v_1)}{\exp(s_1(v_1))}$
			\item<3->$u_1 = \frac{v_1 - t_2(u_2)}{\exp(s_2(u_2))}$
			\item 
			\item<4>\alert{$s_i$ and $t_i$ can be arbitrarily complicated functions:\\
			no need to be invertible themselves,\\ hence can use trainable functions}
			\end{description}
	% \end{center}
	% \begin{proof}
	% 	\begin{enumerate}
	% 	\item<1-| alert@1> Suppose $p$ were the largest prime number.
	% 	\item<2-| alert@2> Let $q$ be the product of the first $p$ numbers.
	% 	\item<3-| alert@3> Then $q + 1$ is not divisible by any of them.
	% 	\item<1-| alert@4> But $q + 1$ is greater than $1$, thus divisible by some prime
	% 	number not in the first $p$ numbers.\qedhere \end{enumerate}
	% 	\end{proof}
	% 	\uncover<5->{The proof used \textit{reductio ad absurdum}.}
\end{frame}


\begin{frame}
	\frametitle{INN Main Building Block: Affine Coupling Layer}
	% \begin{center}
			\begin{description}
			\item<1->[forward process:]\includegraphics[width=9cm]{INN-coupling-layer.png}
			\item<1->[inverse process:]\includegraphics[width=9cm]{INN-coupling-layer-inverse.png}\cite{Ardizzone2_2019}
			% \item<3->\alert{$s_i$ and $t_i$ need not be invertible themselves:\\can use trainable functions}
			\end{description}
	% \end{center}
	% \begin{proof}
	% 	\begin{enumerate}
	% 	\item<1-| alert@1> Suppose $p$ were the largest prime number.
	% 	\item<2-| alert@2> Let $q$ be the product of the first $p$ numbers.
	% 	\item<3-| alert@3> Then $q + 1$ is not divisible by any of them.
	% 	\item<1-| alert@4> But $q + 1$ is greater than $1$, thus divisible by some prime
	% 	number not in the first $p$ numbers.\qedhere \end{enumerate}
	% 	\end{proof}
	% 	\uncover<5->{The proof used \textit{reductio ad absurdum}.}
\end{frame}

\begin{frame}
	\frametitle{Characterization of Invertible Neural Networks}
	\begin{enumerate}
		\item mapping bijective: \alert{has inverse}
		\item forward and inverse mapping \alert{efficiently computable}
		\item forward and inverse mapping with \alert{tractable Jacobian}
	\end{enumerate}
\end{frame}


\begin{frame}
	\frametitle{Training Scheme}
	\begin{center}	
		\begin{tikzpicture}
			\node<1-> [anchor=south west, inner sep=0] (image) at (0,0) {\includegraphics[width=10cm]{INN-training-scheme.png}};
            \draw<1> [red, thick] (7, 1.5) rectangle (10,2.9);
            \draw<1> [red, thick] (7, 0.4) rectangle (10, 1.7);
            \node<1> [align=center, black, font={\Large}] at (5,-0.75) {forward loss: any supervised loss $\mathcal{L}_y(y_i, f_y(x_i))$ \\ + MMD};
            \draw<2> [red, thick] (0.2,0.4) rectangle (2.6,2.9);
            \node<2> [align=center, black, font={\Large}] at (5,-0.75) {backward loss: dual loss $\mathcal{L}_z(q(y, z), p(y)p(z))$};

		\end{tikzpicture}\cite{Ardizzone2_2019}
	\end{center}
\end{frame}



\begin{frame}
	\frametitle{Loss function for backward loss $\mathcal{L}_z$:\\Maximum Mean Discrepancy (MMD)}
	\begin{columns}[T]
		\column{5cm}
			given:\\ $X=\{x_1, \ldots, x_m\} \sim p$, $Y=\{y_1, \ldots, y_n\} \sim q$\\
			\alert{test if $p=q$} 
		\column{5cm}
			\begin{center}
				\includegraphics[width=5cm]{MMD.png}\cite{Smola2006}
			\end{center}
	\end{columns}
	\begin{itemize}
		\item<1->Kullback-Leibler divergence or $L^1$/$L^2$ distance compare $\hat p, \hat q$: \alert{indirect measure}
		\item<2->MMD uses kernel trick: \alert{direct measure}
		\item<2->MMD does \alert{not require explicit Jacobians}
	\end{itemize}
\end{frame}


\section{Examples}

\begin{frame}
	\frametitle{Toy Example: Gaussian Mixture Model}
	\begin{description}
		\item<1->\includegraphics[width=7cm]{toy-truth}\cite{Ardizzone2_2019}
		\item<2->\movie[externalviewer]{\includegraphics[width=7cm]{INN-toy-animation-poster}}{INN-toy-animation.mp4}
	\end{description}
\end{frame}


\begin{frame}
	\frametitle{Real-World Example: Biological Tissue Parameters from Multispectral Image}
	\begin{center}
		\begin{tikzpicture}
			\node<1>{\includegraphics[width=9.25cm]{tissue-data.png}\cite{Ardizzone1_2019}};
			\node<2>{\includegraphics[width=9.25cm]{tissue-table.png}\cite{Ardizzone1_2019}};	
		\end{tikzpicture}
	\end{center}
\end{frame}



\section{Conclusion}


\begin{frame}
	\frametitle{Benefits of the INN Method}
	\begin{itemize}
		\item Very good quantitative and qualitative results
		\item Relatively easy, cheap and straightforward to train
	\end{itemize}

\end{frame}


\begin{frame}
	\frametitle{Challenges of the INN Method}

	\begin{itemize}
		\item How to decide the intrinsic dimension of the data?
		\item How to decide splitting $x$ into $u_1$ and $u_2$?
		\item How to decide permutation of the streams $u_1$ and $u_2$ between coupling layers?
	\end{itemize}
\end{frame}

\begin{frame}
	\frametitle{My Take}
	\begin{center}
		\begin{tikzpicture}
			\node<2->{\includegraphics[width=5cm]{positive.png}};
		\end{tikzpicture}
	\end{center}

\end{frame}


\appendix
\section*{References}
\begin{frame}[allowframebreaks]
\frametitle{References}
	\begin{thebibliography}{10}
		\bibliographystyle{apalike}

		% \beamertemplatebookbibitems
		% % Start with overview books.
	  
		% \bibitem{Author1990}
		%   A.~Author.
		%   \newblock {\em Handbook of Everything}.
		%   \newblock Some Press, 1990.
	   
		  
		\beamertemplatearticlebibitems
		% Followed by interesting articles.
	  
		\bibitem{Ardizzone1_2019}
		  L. Ardizzone et al.
		  \newblock Analyzing Inverse Problems with Invertible Neural Networks.
		  \newblock {\em ICLR 2019 conference paper}, \href{https://arxiv.org/abs/1808.04730}{arXiv:1808.04730},
		  2019.
		
		  \bibitem{Ardizzone2_2019}
		  L. Ardizzone et al.
		  \newblock Analyzing Inverse Problems with Invertible Neural Networks.
		  \newblock {\em Visual Learning Lab Heidelberg}, \href{https://hci.iwr.uni-heidelberg.de/vislearn/inverse-problems-invertible-neural-networks/}{Blog Post},
		  2018.

		  \bibitem{Dinh2016}
		  L. Dinh et al.
		  \newblock Density Estimation using Real NVP.
		  \newblock {\em ICLR 2017 conference paper}, \href{https://arxiv.org/abs/1605.08803}{arXiv:1605.08803},
		  2016.
	
		  \bibitem{Smola2006}
		  A. Smola
		  \newblock Maximum Mean Discrepancy.
		  \newblock {\em ICONIP 2006 conference presentation}, \href{http://alex.smola.org/teaching/iconip2006/iconip_3.pdf}{Alexander Smola's Personal Page (Retrieved Sep 26, 2019)},
		  2006.

		  \bibitem{Mirza2014}
		  M. Mirza and S. Osindero
		  \newblock Conditional Generative Adversarial Nets.
		  \newblock \href{https://arxiv.org/abs/1411.1784}{arXiv:1411.1784},
		  2014. 
		 
		  \bibitem{Blundell2015}
		  C. Blundell et al.
		  \newblock Weight Uncertainty in Neural Networks.
		  \newblock {\em Google DeepMind}, \href{https://arxiv.org/abs/1505.05424}{arXiv:1505.05424},
		  2015.

		  \bibitem{Nix1994}
		  D. Nix and A. Weigend
		  \newblock Estimating the Mean and Variance of the Target Probability Distribution.
		  \newblock {\em ICNN94 conference paper}, \href{https://doi.org/10.1109/ICNN.1994.374138}{DOI:10.1109/ICNN.1994.374138},
		  1994

		  \bibitem{Sunnaker2013}
		  M. Sunn\aa ker et al.
		  \newblock Approximate Bayesian Computation.
		  \newblock {\em PLoS Comput Biol 9(1): e1002803}, \href{https://doi.org/10.1371/journal.pcbi.1002803}{DOI:10.1371/journal.pcbi.1002803},
		  2013

	\end{thebibliography}

\end{frame}
\end{document}
